% SHARD-ID: RC-08G-WEIL-WINDOW-EXTENSION
% SHARD-TITLE: Weil Positivity, Continued: the Windowed Form Is a Compression, the Extension Disc for the Next Trace, and What the Counts Add
% SHARD-SUMMARY: The Connes--Consani--Moscovici chain and its graph analogue never model the traces beyond the window: the window form is the exact compression of the full Weil form, and the model of the higher traces is implicit, the Carath\'eodory--Fej\'er / Pisarenko extension (atoms on the circle plus a noise floor at lag zero).
% SHARD-SUMMARY: Given the traces up to lag K with a positive window, the admissible next trace fills a closed disc (centre the Levinson one-step predictor, radius a ratio of Toeplitz determinants, boundary points the singular windows with K+1 atoms); no estimator from traces and positivity alone beats it; the counts add nonnegativity and a Newton-identity congruence modulo K+1, which pin the next trace of a permutation but not of a q-regular graph, where the disc radius grows by q^{(K+1)/2} in count units; for zeta the mean of the higher traces is the pole term, already exact, so a better higher-trace estimator means more primes.
% SHARD-KEYWORDS: window, compression, Toeplitz, trigonometric moment problem, extension disc, Levinson, Yule-Walker, maximum entropy, Burg, Pisarenko, Caratheodory-Fejer, Schur parameters, Verblunsky, Newton identities, congruence, permutation, Petersen graph, Connes-Consani-Moscovici, zeta spectral triples, higher traces
% SHARD-NOTES: none
% SHARD-SCRIPTS: scripts/weil_window_extension.py

\section{Weil Positivity, Continued: the Window as a Compression and the Extension Disc}
\label{sec:weil-window-extension}

\emph{Question (TJO, 2026-09-19).} The Connes--Consani--Moscovici construction
\cite{connes2025zeta} (read and implemented in the \texttt{zst} and \texttt{ihz} sidequests,
\texttt{notes/zeta-spectral-triples/}) seems to (1) assume a form for the traces $\Tr X^k$ of a
transfer operator, (2) collect the data up to some lag, (3) build a model for the higher traces
from the incomplete data, and (4) use positivity and the Toeplitz structure to read information
off the minimal eigenvector. The characteristic polynomial of $X$ is built from the traces; could
one be cleverer and build a better estimator of $\Tr X^l$ for $l > K$ from the traces up to $K$?
Everything here is elementary and the notebook's own calculation
(\texttt{scripts/weil\_window\_extension.py}, author claude:fable-5.1, unreviewed); the classical
names are given for orientation, not as citations, except the Carath\'eodory--Fej\'er corollary,
which is byte-cited from \cite{connes2025quadratic}.

\subsection{The window is a compression, and the model of the higher traces is implicit}

Throughout, $X$ is a finite transfer operator with trivial multiset $S_{\mathrm{triv}}$ and
critical radius $r$, and $\nuseq{k}$ is its rescaled trace sequence as in
\cref{thm:weil-positivity-finite}, $\nuseq{-k} = \overline{\nuseq{k}}$. For a permutation matrix
$P$ (\cref{sec:graded-permutation}) $r = 1$, $S_{\mathrm{triv}} = \emptyset$ and
$\nuseq{k} = \Tr P^k = \#\mathrm{Fix}(\pi^k)$ (\cref{prop:permutation-rings}); for a
$(q+1)$-regular graph $X = \Hash$, $r = \sqrt q$ and $S_{\mathrm{triv}}$ is the Ihara--Bass
trivial divisor, so that $\nuseq{k} = q^{-k/2}(\Tr\Hash^k - \text{trivial})$.

\begin{definition}[Window matrix, admissible extensions, extremal extensions]
\label{def:window-extension-set}
The \emph{window matrix} of order $K$ is the $(K+1)\times(K+1)$ Hermitian Toeplitz matrix
$\Wform_K = (\nuseq{j-k})_{j,k=0}^{K}$, the matrix of the Weil form $\Wform$ of
\cref{thm:weil-positivity-finite} on test sequences supported in $\{0,\dots,K\}$. The
\emph{admissible set} at order $K$ is the set of values $x\in\mathbb C$ such that the window
matrix of order $K+1$ built from $\nuseq{0},\dots,\nuseq{K},x$ is positive semidefinite. An
\emph{extension} of the data $\nuseq{0},\dots,\nuseq{K}$ is any sequence agreeing with it up to lag
$K$ whose window matrices of every order are positive semidefinite, equivalently (by the
Carath\'eodory--Toeplitz theorem) the Fourier coefficient sequence of a positive measure on the
unit circle. The \emph{Pisarenko extension} is obtained by subtracting the minimal eigenvalue
$\varepsilon_K$ of $\Wform_K$ from the diagonal, taking the $K$ roots $z_1,\dots,z_K$ of the
kernel polynomial $\sum_j\kappa_jz^j$ of $\Wform_K - \varepsilon_K$ and the weights $w_j$ that
reproduce the shifted data, and setting $\nuseq{k} = \sum_j w_jz_j^k + \varepsilon_K\delta_{k,0}$
for all $k$.
\end{definition}

\begin{observation}[The chain never models the higher traces; its implicit model is Pisarenko's]
\label{obs:window-is-compression}
\claimstatus{obs:window-is-compression}{sketched}
In the Connes--Consani--Moscovici chain and in its graph analogue
(\texttt{notes/zeta-spectral-triples/ihara/lanes/theory.md}, IH-3 and IH-11), the window
quadratic form is the exact restriction of the full Weil form to test functions supported in the
window: for a graph, only $\nuseq{0},\dots,\nuseq{K}$ enter $\Wform_K$, and they enter exactly;
for $\zeta$, only the primes $p\le\lambda^2$ enter the form on $[\lambda^{-1},\lambda]$, and the
pole and archimedean terms enter exactly. No value is assigned to the traces beyond the window.
Positivity of every compression is Weil's criterion, and positivity of the compression is
implied by it. The perturbed operator of the chain has as spectrum exactly the roots of the kernel
polynomial of $\Wform_K - \varepsilon_K$ (IH-11(iv)), so the only model of the higher traces the
chain commits to is the Pisarenko extension of \cref{def:window-extension-set}: $K$ atoms on the
circle plus a white-noise floor of size $\varepsilon_K$ at lag zero. This is the minimal-support
extension of the data, and it is the appropriate prior exactly because the true spectral measure
is pure point.
\end{observation}

\subsection{The extension disc}

\begin{citedfact}[Carath\'eodory--Fej\'er corollary]
\label{cit:cf-kernel-circle}
\claimstatus{cit:cf-kernel-circle}{cited}
\sloppy
Connes and van Suijlekom \cite{connes2025quadratic} state and reprove the corollary of the 1911
Carath\'eodory--Fej\'er structure theorem:
{\small\texttt{\detokenize{if $T \in M_n(\mathbb{C})$ is a Hermitian, positive semidefinite Toeplitz matrix of rank $n - 1$, and $\xi \in \ker T$, then the polynomial $P(z) = \sum \xi_j z^j$ has all its zeros on the unit circle.}}}
(\texttt{2511.23257:Araki-final-oct25.tex:752}). In the notation of this section: a Hermitian
positive semidefinite Toeplitz matrix of rank one less than its size has a one-dimensional kernel,
and the polynomial of the kernel vector has all its zeros on the unit circle.
\end{citedfact}

\begin{proposition}[The admissible next trace fills a disc]
\label{prop:extension-disc}
\claimstatus{prop:extension-disc}{sketched}
Let $\Wform_K$ be positive definite and write its last column above the diagonal as
$u(x) = (\bar x,\overline{\nuseq{K}},\dots,\overline{\nuseq{1}})^{\mathsf T}$ for a candidate
$x = \nuseq{K+1}$.
\begin{enumerate}
\item[(i)] The admissible set is the closed disc $|x - c_K|\le r_K$ with
  $c_K = -\overline{(\Wform_K^{-1}w)_0}/(\Wform_K^{-1})_{00}$ (conjugate added 2026-09-24 after the RTP-1 review; all data in this book are real), where $w = u(0)$, and
  $r_K^2 = (\nuseq{0} - w^*\Wform_K^{-1}w)/(\Wform_K^{-1})_{00} + |c_K|^2$.
\item[(ii)] The centre is the one-step linear predictor: with $a$ the solution of the
  Yule--Walker system $\Wform_{K-1}a = -(\nuseq{1},\dots,\nuseq{K})^{\mathsf T}$,
  $c_K = -\sum_{j=1}^{K}a_j\nuseq{K+1-j}$; the radius is the prediction error
  $r_K = \det\Wform_K/\det\Wform_{K-1}$, non-increasing in $K$.
\item[(iii)] $x$ lies on the boundary of the disc iff $\Wform_{K+1}(x)$ is singular, and then its
  kernel is one-dimensional and the kernel polynomial has its $K+1$ roots on the unit circle
  (\cref{cit:cf-kernel-circle}): the boundary points are exactly the extensions by $K+1$ atoms.
\item[(iv)] If the retained measure of $X$ has exactly $K+1$ distinct atoms, the true value
  $\nuseq{K+1}$ lies on the boundary of the disc of order $K$, and every later trace is determined
  by the kernel recurrence of $\Wform_{K+1}$.
\item[(v)] Each further lag adds one free parameter in the closed unit disc (the Schur, or
  Verblunsky, parameter); no estimator of $\nuseq{K+1}$ from $\nuseq{0},\dots,\nuseq{K}$ and
  positivity alone can do better than the disc of (i). The centre $c_K$ is the minimax choice and
  is the extension of maximal entropy (Burg); the Pisarenko extension of
  \cref{def:window-extension-set} lies in the disc.
\end{enumerate}
\end{proposition}

\begin{proof}
(i) By the Schur complement, $\Wform_{K+1}(x)\ge0$ iff
$g(x) := \nuseq{0} - u(x)^*\Wform_K^{-1}u(x)\ge0$. With $u(x) = \bar xe_0 + w$,
$g(x) = \nuseq{0} - |x|^2(\Wform_K^{-1})_{00} - 2\,\mathrm{Re}\big(x\,(\Wform_K^{-1}w)_0\big)
- w^*\Wform_K^{-1}w$, and completing the square gives the disc with the displayed centre and
radius; $r_K^2\ge0$ because the true next trace is admissible. (ii) The Yule--Walker system is the
normal equation for the best monic predictor in $L^2$ of the measure whose Fourier coefficients
are the $\nuseq{k}$, and $g(c_K) = \max g$ is its residual; $(\Wform_K^{-1})_{00} =
\det\Wform_{K-1}/\det\Wform_K$ by Cramer's rule and Toeplitz shift invariance, which turns
$g(c_K)/(\Wform_K^{-1})_{00}$ into $(\det\Wform_K/\det\Wform_{K-1})^2$; the script checks both
identities to $10^{-9}$ on every window of the two examples rather than carrying the algebra
here, which is why the status is \texttt{sketched}. Monotonicity of $r_K$ is the monotonicity of
prediction errors under refinement. (iii) $\det\Wform_{K+1}(x) = g(x)\det\Wform_K$, so
singularity is $g = 0$, the boundary; the kernel is one-dimensional because $\Wform_K$ is
definite; \cref{cit:cf-kernel-circle} gives the circle. (iv) With $m = K+1$ distinct atoms the
window matrices have rank at most $m$, so $\Wform_{K+1}$ (size $K+2$) is singular while
$\Wform_K$ (size $K+1$) is definite (a Vandermonde argument at $m$ distinct nodes); singular
means boundary by (iii), and the kernel vector of $\Wform_{K+1}$ is the polynomial vanishing on
the atoms, whose recurrence propagates every later trace. (v) is the Schur parametrisation of the
truncated trigonometric moment problem (Verblunsky's theorem), stated here for orientation and
not used by anything else in this section.
\end{proof}

\subsection{What the counts add}

The traces of a transfer operator built from letters are counts, and counts carry information
that positivity does not see.

\begin{proposition}[Newton congruence]
\label{prop:newton-congruence}
\claimstatus{prop:newton-congruence}{sketched}
Let $X$ be a square matrix with integer entries and $N_k = \Tr X^k$. Then for every $K\ge0$,
$N_{K+1}\equiv -\sum_{i=1}^{K}c_iN_{K+1-i}\pmod{K+1}$, where $\det(1-uX) = \sum_ic_iu^i$ and the
$c_i$ are integers determined by $N_1,\dots,N_i$. So the traces up to lag $K$ determine the next
trace modulo $K+1$.
\end{proposition}

\begin{proof}
Newton's identity for the characteristic polynomial reads
$(K+1)c_{K+1} = -\sum_{i=0}^{K}c_iN_{K+1-i}$ with $c_0 = 1$; the $c_i$ are integers because
$\det(1-uX)\in\mathbb Z[u]$. Reduce modulo $K+1$.
\end{proof}

\begin{observation}[Rescaling defeats integrality; for $\zeta$ the mean of the higher traces is the
pole]
\label{obs:rescaling-and-zeta}
\claimstatus{obs:rescaling-and-zeta}{sketched}
(i) For a permutation ($r = 1$) the disc of \cref{prop:extension-disc}, the half-line
$N_{K+1}\ge0$ and the congruence of \cref{prop:newton-congruence} act on the same integer, and
can pin it before the window is exact (\cref{num:weil-window-extension}: pinned at $K = 3$ for
$(1\,2)(3\,4\,5)$, one lag before the critical window). (ii) For a $(q+1)$-regular graph the disc
lives in the rescaled variable and has radius $q^{(K+1)/2}r_K$ in count units, which exceeds one
on every positive-definite window of the Petersen graph (radii $35.5$, $40.6$, $49.4$ in counts
at $K = 1,2,3$), so integrality pins nothing; only the Ramanujan-side information (positivity) and
the exact critical window remain. (iii) For $\zeta$ the traces are $\vM(n)$. The mean of the
traces beyond the window is the prime number theorem, which is exactly the pole term of the
explicit formula and enters the window form exactly; the higher powers $p^m > \lambda^2$ of the
known primes lie outside the support of the windowed form and are invisible to it; what is
unknown beyond the window is the fluctuation $\vM(n) - 1$, which by the explicit formula is
the zeros. A better estimator of the higher traces of $\zeta$ therefore means more primes, i.e.\ a
larger window. The freedom that remains is the choice of extremal extension in
\cref{prop:extension-disc}(v), and the content of \cite{connes2025zeta} is the rate at which the
Pisarenko atoms converge to the true zeros as the window grows (about $5.4$ digits per unit of
$\lambda^2$ in the \texttt{zst} benchmark), the prolate step that a finite divisor cannot
exercise because its disc radius reaches zero at the critical window
(\cref{prop:extension-disc}(iv)).
\end{observation}

\begin{numerical}[Windows, discs, congruences: the permutation $(1\,2)(3\,4\,5)$ and the Petersen graph]
\label{num:weil-window-extension}
\claimstatus{num:weil-window-extension}{numerical}
\texttt{scripts/weil\_window\_extension.py}, $45$ checks. Permutation: $\Tr P^k =
5,0,2,3,2,0,5,\dots$; $\Wform_K$ definite for $K\le3$, singular from $K = 4$ (four distinct
atoms); at $K = 1,2,3$ the disc centre equals the Levinson predictor and the radius equals
$\det\Wform_K/\det\Wform_{K-1}$ ($5$, $4.2$, $2.0571$), positivity of $\Wform_{K+1}(x)$ holds
exactly on the disc ($61$ sample points per window), the truth lies inside, both boundary points
give singular windows with all kernel roots on the circle, the Pisarenko extension lies inside,
the Newton congruence holds, and the integers $\ge0$ in the disc with the right residue are
$\{0,2,4\}$, $\{0,3\}$, $\{2\}$: pinned at $K = 3$, where the truth $2$ sits on the boundary and
Pisarenko alone predicts $-0.5$ (repeated atom at $1$). At $K = 4,5$ the kernel recurrence returns
$\nuseq{5} = 0$, $\nuseq{6} = 5$ exactly and the kernel roots are the four distinct eigenvalues.
Petersen ($q = 2$, $18$ retained eigenvalues at angles $\pm\arccos(1/2\sqrt2)$ and $\pm3\pi/4$):
rescaled traces reproduce $\Tr\Hash^k = 30,0,0,0,0,120,120,0$ exactly; $\Wform_K$ definite for
$K\le3$; the truth lies in every disc and on the boundary at $K = 3$; the congruence holds; the
counts-unit radii are $35.5$, $40.6$, $49.4$ with $26$, $19$, $25$ admissible integers.
\end{numerical}

\paragraph{Not established, and one computation proposed.}
None of this is a statement about $\zeta$ beyond \cref{obs:rescaling-and-zeta}(iii), which is an
argument, not a theorem. The disc radius $r_K$ is a rigorous measure of ignorance about the next
lag; in \texttt{ihz} it is one determinant ratio per window, and combined with the congruence
and nonnegativity it would show how far below the critical window a graph divisor is already
forced. For \texttt{zst} the analogous quantity, the Schur-complement radius of the window
Loewner form as $\lambda$ grows, would bound rigorously what the primes beyond $\lambda^2$ can do
to the form; comparing its decay in $\lambda$ with the observed $e^{-4\pi\lambda^2}$ convergence
of the first zero would say whether that convergence is explained by the shrinking envelope or by
something specific to the arithmetic. Not run.
